\chapter*{Exposé XXV\\The Existence Theorem}
\addcontentsline{toc}{chapter}{Exposé XXV. The Existence Theorem}
\begin{center}
\textit{by M. Demazure}
\end{center}

\footnote{Editors' note: Version of 13 October 2024.}
For completeness, in this Exposé we give a proof of the existence theorem for split
groups. Like Chevalley's original proof (``Sur certains schémas de groupes
semi-simples'', \textit{Séminaire Bourbaki}, May 1961, no.~219), it rests on the
existence of complex semisimple algebraic groups of every possible type. Cartier
gave the principle of a more satisfactory proof, establishing directly the
existence of a simply connected split semisimple \(\mathbb Z\)-group associated
with a given Cartan matrix (unpublished).\footnote{Editors' note: Such a principle
was also sketched by B.~Kostant \emph{[Ko66]}; a complete proof, using quantum
groups, was given more recently by G.~Lusztig \emph{[Lu09]}.}
Let us point out, however, that the difficulty is not to give an explicit
construction of a group scheme, but to verify that the group so constructed
really satisfies the required conditions---essentially, that its fibers are
indeed smooth and reductive.

\SGASection{1}{Statement of the theorem}
\SGARefTarget{sga3:xxv:section:1}

\begin{sgastatement}{Theorem 1.1}
Let \(S\) be a nonempty scheme. The functor
\[
  G\longmapsto R(G)
\]
is an equivalence from the category of pinned reductive \(S\)-groups to the
category of pinned reduced root data.
\end{sgastatement}
\SGARefTarget{sga3:xxv:theorem:1:1}

By the uniqueness theorem (Exposé XXIII, 4.1), the preceding statement is
equivalent to the following.

\begin{sgastatement}{\SGAInlineRefTarget{sga3:xxv:corollary:1-2}Corollary 1.2 (Existence of split groups)}
For every reduced root datum \(R\), there exists a split reductive
\(\mathbb Z\)-group \(G\) such that \(R(G)\simeq R\).\footnote{Editors' note:
The work \emph{[BT84]} of F.~Bruhat and J.~Tits also contains a variant of
Chevalley's construction (\emph{loc.\ cit.}, 2.2.3--2.2.5 and \S3.2), which
in particular gives a smooth affine \(\mathbb Z\)-group \(G\) with connected
fibers, possessing a split maximal torus, and whose generic fiber is a reductive
\(\mathbb Q\)-group of type \(R\). The reductivity of the geometric fibers of
\(G\) follows from the study of the unipotent radical of a special fiber in
\emph{loc.\ cit.}, 4.6.12 and 4.6.15 (valid for more general \(G\), associated
with a valued root datum); but it is simpler to deduce it from the description
of \(\mathfrak g=\operatorname{Lie}(G)\), from Exposé XIX, 1.12(iii), and from
the existence of the elements \(w_\alpha(X)\) of \SGARefLink{sga3:xx:3.1:item:iv}{Exposé XX, 3.1(iv)}; see
\emph{[BT84]}, 3.2.1.}
\end{sgastatement}
\SGARefTarget{sga3:xxv:corollary:1:2}

In particular:

\begin{sgastatement}{Corollary 1.3}
Let \(k\) be a field. For every splittable reductive \(k\)-group \(G_k\), there
exists a reductive \(\mathbb Z\)-group \(G\) such that
\[
  G\otimes_{\mathbb Z}k\simeq G_k.
\]
\end{sgastatement}
\SGARefTarget{sga3:xxv:corollary:1:3}

Conversely, first observe that to prove \SGARefLink{sga3:xxv:corollary:1-2}{1.2} it is enough, by \SGARefLink{sga3:xxii:prop:4.3.1}{Exposé XXII, 4.3.1},
and \SGARefLink{sga3:XXI:6.5.10}{Exposé XXI, 6.5.10}, to consider the case in which the root datum \(R\) is
simply connected (and even irreducible, if desired, by \SGARefLink{sga3:XXI:7.1.6}{Exposé XXI, 7.1.6}).

On the other hand, under the hypotheses of \SGARefLink{sga3:xxv:corollary:1:3}{1.3} the group scheme \(G\) has
constant type (since \(\operatorname{Spec}(\mathbb Z)\) is connected), hence type
\(R(G_k)\). By \SGARefLink{sga3:XXIII:5.9}{Exposé XXIII, 5.9}, the validity of \SGARefLink{sga3:xxv:corollary:1:3}{1.3} for a given group \(G_k\)
therefore entails the existence of a split \(\mathbb Z\)-group of type
\(R(G_k)\).\footnote{Editors' note: In fact, \SGARefLink{sga3:XXIII:5.9}{Exposé XXIII, 5.9} is unnecessary:
the present Exposé constructs, for every semisimple \(\mathbb Q\)-group
\(G_{\mathbb Q}\), a semisimple \(\mathbb Z\)-group \(G\) of the same type as
\(G_{\mathbb Q}\), equipped with a split maximal torus \(T\); hence \(G\) is
split by \SGARefLink{sga3:xxii:prop:2.2}{Exposé XXII, 2.2}.}

Thus, to prove \SGARefLink{sga3:xxv:corollary:1-2}{1.2} and hence \SGARefLink{sga3:xxv:theorem:1:1}{1.1}, it is enough to prove \SGARefLink{sga3:xxv:corollary:1:3}{1.3} when \(k\) has
characteristic zero (for example, \(k=\mathbb C\)) and \(G_k\) is simply
connected (and in particular semisimple), together with the following
proposition.

\begin{sgastatement}{Proposition 1.4}
For every simply connected reduced root datum \(R\), there exists a semisimple
algebraic \(\mathbb C\)-group of type \(R\).
\end{sgastatement}
\SGARefTarget{sga3:xxv:proposition:1:4}

One may prove \SGARefLink{sga3:xxv:proposition:1:4}{1.4} as follows. First, there exists a complex semisimple Lie
algebra \(\mathfrak g\) of type \(R\); see, for example, N.~Jacobson,
\textit{Lie Algebras}, Chapter VII, Theorem 5.\footnote{Editors' note:
Let \(\Delta\) be a base of \(R\), and let \(\widetilde{\mathfrak g}\) be the
complex Lie algebra generated by
\((e_\alpha,f_\alpha,h_\alpha)_{\alpha\in\Delta}\), subject to the relations
\([h_\alpha,h_\beta]=0\), \([e_\alpha,f_\beta]=h_\alpha\) if
\(\beta=\alpha\) and \(=0\) otherwise,
\([h_\alpha,e_\beta]=(\alpha^\ast,\beta)e_\beta\), and
\([h_\alpha,f_\beta]=-(\alpha^\ast,\beta)f_\beta\).
In \emph{loc.\ cit.}, \(\mathfrak g\) is defined as the quotient of
\(\widetilde{\mathfrak g}\) by the intersection of the kernels of its
finite-dimensional irreducible representations. In \emph{[Se66]},
\S VI.5, Theorem 9 (see also \emph{[BLie]}, VIII, \S4.3, \SGARefLink{sga3:xxv:section:1}{Theorem 1}), it is
shown that \(\mathfrak g\) is the quotient of \(\widetilde{\mathfrak g}\) by
the relations
\[
 \operatorname{ad}(e_\alpha)^{1-(\alpha^\ast,\beta)}(e_\beta)=0,\qquad
 \operatorname{ad}(f_\alpha)^{1-(\alpha^\ast,\beta)}(f_\beta)=0.
\]
For an explicit description of the structure constants, and in particular
the choice of signs, see \SGARefLink{sga3:XXIII:6.5}{Exposé XXIII, 6.5} and 6.7, as well as
\emph{[Ti66]}, \S4, \SGARefLink{sga3:xxv:section:1}{Theorem 1}, and, for types \(A,D,E\),
\emph{[Sp98]}, 10.2.5.}
Then \(G=\operatorname{Aut}(\mathfrak g)^0\) is a semisimple algebraic
\(\mathbb C\)-group of type \(\operatorname{ad}(R)\).\footnote{Editors' note:
One may suppose \(R\) irreducible, hence \(\mathfrak g\) simple. By a general
argument, \(\operatorname{Lie}(G)\) is the complex Lie algebra of derivations
of \(\mathfrak g\); see \emph{[DG70]}, \S II.4, Proposition 2.3. These
derivations are all inner; see \emph{[BLie]}, I, \S6.1, Corollary 3 to
\SGARefLink{sga3:xxv:section:1}{Proposition 1}. Thus \(\operatorname{Lie}(G)=\mathfrak g\), so \(G\) has no
invariant subgroup of positive dimension and is semisimple. Its root system
is consequently the same as that of \(\mathfrak g\). Moreover, the center of
\(G\) acts both trivially and faithfully on \(\mathfrak g\), and is therefore
trivial; hence \(G\) is adjoint.}
By \textit{Bible}, \S23.1, \SGARefLink{sga3:xxv:section:1}{Proposition 1}, it follows that there exists a
semisimple \(\mathbb C\)-group of type \(R\).

The remainder of this Exposé is devoted to the proof of \SGARefLink{sga3:xxv:corollary:1:3}{1.3} when \(k\) has
characteristic zero and \(G_k\) is semisimple. The proof proceeds in two
stages: the construction of a ``piece of \(\mathbb Z\)-group scheme''
(\S2), and the study of the group obtained by applying ``Weil's theorem''
(\S3).

To avoid confusion, in \S2 we shall not use the abbreviated notation \(X_k\)
for the \(k\)-scheme \(X\otimes_{\mathbb Z}k\), where \(X\) is a
\(\mathbb Z\)-scheme.

